%\documentclass[12pt,leqno]{amsart}
\renewcommand{\baselinestretch}{1.1}
\usepackage[utf8]{inputenc}
\usepackage{fullpage}
\usepackage[colorlinks,linkcolor=blue,citecolor=blue]{hyperref}
\usepackage{graphicx}
\usepackage{amsmath}
\usepackage{amsthm}
\usepackage{amsfonts}
\usepackage{amssymb}
\usepackage{bbm}
\usepackage{mathtools}
\usepackage{extarrows}
\usepackage{array}   % for \newcolumntype macro
\newcolumntype{L}{>{$}l<{$}}

\usepackage{tikz-cd}
\usepackage{pgf,tikz}
\usetikzlibrary{
  knots,
  hobby,
  decorations.pathreplacing,
  shapes.geometric,
  calc
}

%\tikzset{
% knot diagram/every strand/.append style={
%    ultra thick,
%    red
%  },
% show curve controls/.style={
%    postaction=decorate,
%    decoration={show path construction,
%      curveto code={
%        \draw [blue, dashed]
%        (\tikzinputsegmentfirst) -- (\tikzinputsegmentsupporta)
%        node [at end, draw, solid, red, inner sep=2pt]{};
%        \draw [blue, dashed]
%        (\tikzinputsegmentsupportb) -- (\tikzinputsegmentlast)
%        node [at start, draw, solid, red, inner sep=2pt]{}
%        node [at end, fill, blue, ellipse, inner sep=2pt]{}
 %       ;
 %     }
 %   }
 % },
  %show curve endpoints/.style={
  %  postaction=decorate,
  %  decoration={show path construction,
  %    curveto code={
  %      \node [fill, blue, ellipse, inner sep=2pt] at (\tikzinputsegmentlast) {}
   %     ;
   %   }
  %  }
  %}
%}

\usepackage{mathrsfs}
\usetikzlibrary{arrows}

%\usepackage{amscd,epsf, latexsym}
%\usepackage[dvips]{graphicx}
%%%%%%%%%%%%%%%%%%%%%%%%%%%%%%%%%%%%%%%%%%%%%%%%%%%%%%%
% symbol macros %
%--------Theorem Environments--------
%\theoremstyle{plain} --- default
\newtheorem{Thm}{Theorem}[section]
\newtheorem{thm}{Theorem}[subsection]
\newtheorem{Cor}[Thm]{Corollary}
\newtheorem{prop}[Thm]{Proposition}
\newtheorem{lem}[Thm]{Lemma}
\newtheorem{conj}[Thm]{Conjecture}
\newtheorem{quest}[thm]{Question}
\newtheorem{defn}[Thm]{Definition}
\newtheorem{Fact}{Fact}

\theoremstyle{definition}
\newtheorem{defns}[thm]{Definitions}
\newtheorem{con}[thm]{Construction}
\newtheorem{exmp}[Thm]{Example}
\newtheorem{exmps}[thm]{Examples}
\newtheorem{notn}{Notation}
\newtheorem{notns}[thm]{Notations}
\newtheorem{addm}[thm]{Addendum}
\newtheorem{exer}[thm]{Exercise}


\theoremstyle{remark}
\newtheorem{ques}[Thm]{Question}
\newtheorem{rem}[thm]{Remark}
\newtheorem{rems}[thm]{Remarks}
\newtheorem{warn}[thm]{Warning}
\newtheorem{sch}[thm]{Scholium}
\newtheorem{rmk}{Remark}[section]
\newcommand{\one}{\mathbf{1}}
\newcommand{\mcL}{\mathcal{L}}
\usepackage{eucal}
\usepackage{graphicx}
\usepackage{multirow}
\usepackage[all,knot]{xy}
\newcommand{\er}[1]{{\textcolor{blue}{#1}}}
\newcommand{\pg}[1]{{\textcolor{red}{#1}}}
\DeclareMathOperator{\Rep}{Rep}
\DeclareMathOperator{\Stab}{Stab}
\DeclareMathOperator{\Ising}{Ising}
\DeclareMathOperator{\sVec}{sVec}
\DeclareMathOperator{\Vect}{Vec}
\DeclareMathOperator{\Sem}{Sem}
\DeclareMathOperator{\Mod}{Mod}
\DeclareMathOperator{\Out}{Out}
\DeclareMathOperator{\PGL}{PGL}
\DeclareMathOperator{\PSL}{PSL}
\DeclareMathOperator{\SL}{SL}
\DeclareMathOperator{\GL}{GL}
\DeclareMathOperator{\SMod}{SMod}
\DeclareMathOperator{\LMod}{LMod}
\DeclareMathOperator{\PMod}{PMod}
\DeclareMathOperator{\End}{End}
\DeclareMathOperator{\Add}{Add}
\DeclareMathOperator{\Gal}{Gal}
\DeclareMathOperator{\Tr}{Tr}
\DeclareMathOperator{\Fun}{Fun}
\DeclareMathOperator{\BrPic}{BrPic}
\DeclareMathOperator{\Breq}{Breq}
\DeclareMathOperator{\Pic}{Pic}
\DeclareMathOperator{\Irr}{Irr}
\DeclareMathOperator{\id}{id}
\DeclareMathOperator{\Id}{Id}
\DeclareMathOperator{\obj}{obj}
\DeclareMathOperator{\Obj}{Obj}
\headheight=5pt \headsep=18pt
\footskip=18pt
\textheight=47pc \topskip=10pt
\textwidth=37pc
\calclayout
\newcommand{\RNum}[1]{\uppercase\expandafter{\romannumeral #1\relax}}
\newcommand{\mcB}{\mathcal{B}}
\newcommand{\mcC}{\mathcal{C}}
\newcommand{\mcD}{\mathcal{D}}
\newcommand{\BR}{\mathbb{R}}
\newcommand{\BZ}{\mathbb{Z}}
\newcommand{\Zn}[1]{\BZ/{#1}\BZ}
\DeclareMathOperator{\Aut}{Aut}
\DeclareMathOperator{\Isom}{Isom}

\newcommand{\tY}{\tilde{Y}}
\newcommand{\tX}{\tilde{X}}
\newcommand{\pz}{(0)}
\newcommand{\po}{(1)}
\newcommand{\Tet}[6]{\Big\langle\begin{matrix}
  #1 & #2 & #3 \\
  #4 & #5 & #6 
\end{matrix}\Big\rangle}
\newcommand{\Sym}[6]{\Big\{\begin{matrix}
  #1 & #2 & #3 \\
  #4 & #5 & #6 
\end{matrix}\Big\}}

\usepackage{graphicx}
\graphicspath{ {./graph/} }
\renewcommand{\bot}{\boxtimes}

\newcommand{\fg}{\mathfrak{g}}
\newcommand{\fG}{\mathfrak{G}}
\newcommand{\fS}{\mathfrak{S}}
\newcommand{\fs}{\mathfrak{s}}
\newcommand{\ft}{\mathfrak{t}}
\newcommand{\fa}{\mathfrak{a}}
\newcommand{\fd}{\mathfrak{d}}
\newcommand{\fE}{\mathfrak{E}}
\newcommand{\so}{\mathfrak{so}}
\newcommand{\fm}{\mathfrak{m}}
\newcommand{\fF}{\mathfrak{F}}

\newcommand{\PP}{\mathcal{P}}
\newcommand{\BB}{\mathcal{B}}
\newcommand{\BC}{\mathbb{C}}
\newcommand{\BN}{\mathbb{N}}
\newcommand{\BQ}{\mathbb{Q}}
\newcommand{\BF}{\mathbb{F}}
\renewcommand{\AA}{\mathcal{A}}
\newcommand{\CC}{{\mathcal{C}}}
\newcommand{\DD}{\mathcal{D}}
\newcommand{\II}{\mathcal{I}}
\newcommand{\HH}{\mathcal{H}}
\newcommand{\LL}{\mathcal{L}}
\newcommand{\ZZ}{\mathcal{Z}}
\newcommand{\EE}{\mathcal{E}}
\def\FF{\mathcal{F}}
\newcommand{\RR}{\mathcal{R}}
\newcommand{\WW}{\mathcal{W}}
\newcommand{\VV}{\mathcal{V}}
\def\SS{\mathcal{S}}
\def\MM{\mathcal{M}}
\def\TT{\mathcal{T}}
\def\doubleunderline#1{\underline{\underline{#1}}}

\newcommand{\red}[1]{\textcolor{red}{#1}}
\newcommand{\blue}[1]{\textcolor{cyan}{#1}}
\newcommand{\IH}{\underline{\operatorname{Hom}}}
\newcommand{\Hom}{\operatorname{Hom}}
\newcommand{\bt}{\boxtimes}
\newcommand{\Znx}[1]{(\BZ/{#1}\BZ)^{\times}}
\newcommand{\Gn}[1]{\gal(\BQ_{#1})}
